\documentclass[twocolumn]{revtex4-2}

\usepackage{amsmath,amssymb}
\usepackage{graphicx}
\usepackage{booktabs}
\usepackage{xcolor}
\usepackage{hyperref}
\usepackage{siunitx}

\newcommand{\Mpl}{\bar M_{\rm Pl}}
\newcommand{\Mphi}{M_\phi}
\newcommand{\fzero}{f_0}
\newcommand{\Vrho}{V_{\rm vest}}
\newcommand{\nS}{n_s}
\newcommand{\rT}{r}
\newcommand{\Ne}{N_e}
\newcommand{\eps}{\varepsilon}

% Auto-generated by update_counts.py — 2026-04-24T19:44:03.140400
% Do not edit manually. Run: uv run python scripts/update_counts.py
\newcommand{\totaltheorems}{3993}
\newcommand{\substantivetheorems}{3882}
\newcommand{\placeholdertheorems}{111}
\newcommand{\axiomcount}{1}
\newcommand{\sorrycount}{0}
\newcommand{\sorrypercent}{0.0}
\newcommand{\leanmodules}{167}
\newcommand{\leandefinitions}{3297}
\newcommand{\pythonmodules}{68}
\newcommand{\testfiles}{59}
\newcommand{\totaltests}{2836}
\newcommand{\figurecount}{111}
\newcommand{\notebookcount}{54}
\newcommand{\papercount}{19}
\newcommand{\aristotleproved}{322}
\newcommand{\aristotleruns}{44}
% --- Paper 18 / Phase 5t per-section counts (FermiHubbardDimer.lean) ---
\newcommand{\fhdTotal}{143}
\newcommand{\fhdWaveTwo}{13}
\newcommand{\fhdWaveThree}{6}
\newcommand{\fhdWaveFour}{22}
\newcommand{\fhdWaveFive}{22}
\newcommand{\fhdWaveSixABC}{38}
\newcommand{\fhdWaveSixExtra}{15}
\newcommand{\fhdWaveSeven}{19}
\newcommand{\fhdWaveEight}{8}


\begin{document}

\title{Vestigial natural-branch slow-roll inflation cannot reproduce
Planck/BICEP $(n_s,\,r)$: a structural no-go from the geometric
$\eta$-problem and graceful-exit failure}

\author{John G. Roehm}
\affiliation{SK-EFT Hawking Project (Phase 6b Wave 3)}

\date{\today}

\begin{abstract}
We test whether vestigial-phase coupling dynamics (the
\texttt{VestigialEOS.rho\_vest} potential, anchored in the project's
Phase 5y closure) can support slow-roll inflation compatible with
Planck~2018 and BICEP/Keck~2021. The natural-branch potential
$\Vrho(\tau) = \fzero (1-\tau^2)(5\tau^2-1)$ has a
positive-energy hilltop at $\tau_{\max}^2 = 3/5$ where
$\Vrho(\tau_{\max}) = 4\fzero/5$ and $V''_{\tau\tau}(\tau_{\max}) = -24\fzero$,
producing a closed-form slow-roll-$\eta$ that is
\emph{independent of $\fzero$}: $\eta_{\rm hilltop} = -30\,(\Mpl/\Mphi)^2$.
Slow-roll viability $|\eta|<1$ therefore requires
$\Mphi > \sqrt{30}\,\Mpl \approx 5.48\,\Mpl$ — strictly super-Planckian
and outside the EFT validity window. The Planck $2\sigma$ window on $\nS$
tightens this to $\Mphi > 36\,\Mpl$. Even at the formally compatible
$\Mphi \simeq 41\,\Mpl$, the canonical 50--65 e-fold window is overshot;
graceful exit fails. The structural no-go ships in Lean~4.29 with $15$
substantive theorems plus $1$ summary marker, $0$ sorry, $0$ new axioms
(verified \texttt{propext, Classical.choice, Quot.sound} only), plus a
Python mirror with $48$ pytest cross-layer-parity cases. The result joins
Phase~5y (vestigial-DE no-go vs.\ DESI) and Phase~6b W2 (CMB-$\ell$
gradient-instability) in a coherent project-wide ledger of
vestigial natural-branch falsifications.
\end{abstract}

\maketitle

\section{Introduction}

The post-SK-EFT roadmap's most ambitious cosmology target was a
slow-roll inflation derivation grounded in the vestigial-phase
order-parameter dynamics already formalized at the dark-energy and
perturbation-theory levels in Phase~5y and Phase~6b W2. The roadmap
flagged this wave as ``highest-risk in Phase 6b'' \cite{Phase6bRoadmap},
with cleanly publishable falsifiable status either direction:
quantitative $(\nS, \rT)$ matches to Planck/BICEP would be substantial,
and a structural no-go would parallel the Phase~5y dark-energy no-go.

This paper documents the structural no-go. The vestigial natural-branch
potential's geometry forbids Planck-compatible slow-roll inflation
within the EFT validity window — and the obstruction is geometric, not
parametric, and survives any reasonable $\fzero$ rescaling.

The result joins the project's accumulating ledger of vestigial-branch
falsifications:
\begin{itemize}
\item Phase~5y: vestigial natural-branch dark energy is incompatible
   with DESI Year-1 constraints
   (\texttt{VestigialEOS.vestigial\_not\_in\_desi\_region});
\item Phase~6b W2: CMB-$\ell$ admissibility fails for the vestigial
   gradient-instability regime ($c_s^2 = -1/3$ at $\tau=0$,
   \texttt{CosmologicalPerturbations.vestigial\_at\_zero\_falsifies\_H\_StableSpectrum\_via\_amplitude});
\item This paper: vestigial natural-branch slow-roll inflation cannot
   reproduce $(\nS, \rT)$ at sub-Planckian $\Mphi$, and overshoots
   $\Ne$ at any $\Mphi$ where it could.
\end{itemize}

\section{Setup}
\label{sec:setup}

\paragraph{Vestigial-phase potential.}
Let $\tau\in[0,1]$ be the vestigial-phase order parameter (Phase~5y
\texttt{VestigialEOS}; $\tau=0$ deep vestigial, $\tau=1$ melted).
The Phase~5y closed form is
\begin{align}
\Vrho(\tau) &= \rho_{\rm vest}(\tau)
   = \fzero\,(1-\tau^2)(5\tau^2 - 1).
\end{align}
The potential vanishes at $\tau = 1/\sqrt{5}$ and $\tau = 1$, has an
anti-de-Sitter minimum at $\tau=0$ ($V=-\fzero$), and a positive-energy
hilltop at $\tau_{\max} = \sqrt{3/5}$ with $\Vrho(\tau_{\max}) = 4\fzero/5$.
Inflation requires $V>0$, restricting $\tau_*\in(1/\sqrt{5},\,1)$.

\paragraph{Slow-roll formalism.}
Canonical inflaton $\phi = \Mphi\,\tau$. In reduced-Planck convention
$(\Mpl)$ the slow-roll parameters are
\begin{align}
\eps &= \tfrac{1}{2}\Mpl^2\,(V_\phi/V)^2,
& \eta &= \Mpl^2\,V_{\phi\phi}/V, \\
\nS &= 1 - 6\eps + 2\eta,
& \rT &= 16\eps,
\end{align}
with $V_\phi = (1/\Mphi)V_\tau$ and $V_{\phi\phi}=(1/\Mphi^2)V_{\tau\tau}$.
At the hilltop $V_\tau = 0$, hence $\eps_{\rm hilltop}=0$ and
$\nS = 1 + 2\eta$ exactly.

\paragraph{Microscopic anchor.}
The vestigial free-energy depth $\fzero$ is fixed by the GL ansatz
$\fzero = \pi N_*\,T_{c,{\rm vest}}^2 / 4K(\kappa_*)$ (Phase~5y H4
EQ.111). The inflaton kinetic mass $\Mphi$ is in principle set by the
microscopic kinetic term; for the natural-branch analysis we treat
$\Mphi$ as a free dimensional parameter to be constrained by Planck/BICEP.

\section{The $\eta$-problem at the natural hilltop}
\label{sec:eta}

The hilltop is where slow-roll naturally lives because $\eps=0$ there.
The substantive constraint is on $\eta$. From
$V_{\tau\tau}(\tau_{\max}) = -24\fzero$ and $V(\tau_{\max}) = 4\fzero/5$:
\begin{align}
\eta_{\rm hilltop}
   &= \left(\frac{\Mpl}{\Mphi}\right)^{\!2}\,
      \frac{V_{\tau\tau}(\tau_{\max})}{V(\tau_{\max})}
   = -30\,\left(\frac{\Mpl}{\Mphi}\right)^{\!2}.
\label{eq:eta-closed-form}
\end{align}
The $\fzero$-dependence cancels exactly. The $\eta$-problem is therefore
\emph{geometric}, set by the $V''/V$ ratio of the vestigial potential at
its own positive-energy maximum, not by the energy depth.

\paragraph{Slow-roll viability $|\eta|<1$.} Equation~\eqref{eq:eta-closed-form}
implies
\begin{equation}
|\eta_{\rm hilltop}| < 1
\;\;\Longleftrightarrow\;\;
\Mphi > \sqrt{30}\,\Mpl \approx 5.48\,\Mpl.
\end{equation}
This is the minimal super-Planckian scale that allows the slow-roll
expansion to make sense at all. Lean theorem
\texttt{etaAtHilltop\_abs\_lt\_one\_implies\_super\_Planckian}
formalizes this: $|\eta|<1 \Longrightarrow \Mphi^2 > 30\,\Mpl^2$.

\paragraph{Planck-$2\sigma$ tightening.} At the hilltop $\nS = 1 + 2\eta$,
so $\nS \in [0.9565, 0.9733]$ (Planck~2018 $2\sigma$) translates to
\begin{equation}
\eta_{\rm hilltop} \in [-0.0218, -0.0134],
\end{equation}
which by~\eqref{eq:eta-closed-form} forces
\begin{equation}
\Mphi \in [\,37.1\,\Mpl,\;47.4\,\Mpl\,].
\label{eq:Mphi-window}
\end{equation}
The full window is super-Planckian; the central Planck value
$\nS = 0.9649$ corresponds to $\Mphi \simeq 41.4\,\Mpl$.
Lean theorem \texttt{nSAtHilltop\_planck\_compatible\_implies\_super\_Planckian}
ships the safe under-estimate $\Mphi^2 > 1300\,\Mpl^2$ (i.e.,
$\Mphi > 36\,\Mpl$).

\paragraph{Sub-Planckian falsification.} At $\Mphi = \Mpl$ (the canonical
sub-Planckian choice), $\eta_{\rm hilltop} = -30$ and
\begin{equation}
\nS\big|_{\Mphi=\Mpl} = 1 + 2\cdot(-30) = -59,
\end{equation}
deviating from Planck by $|\nS - 0.9649| = 59.96$ — wildly off-scale, not
a borderline failure. Lean theorem
\texttt{vestigial\_natural\_branch\_inflation\_falsified} formalizes
this: $\Mphi \le \Mpl \Longrightarrow \nS_{\rm hilltop} < 0.9565$.

\section{The graceful-exit (e-fold) overshoot}
\label{sec:efolds}

Even in the super-Planckian window~\eqref{eq:Mphi-window} where $(\nS,\rT)$
become formally Planck-compatible, the second slow-roll observable
$\Ne$ overshoots the canonical $[50, 65]$ window. We confirmed this
numerically with a 2574-point scan over $(f_0, \Mphi, \tau_*)$:
every $(\nS,\rT)$-passing point has $\Ne \ge 71.8$, with median
$\Ne \simeq 319$ across the passing subset. The minimum-$\Ne$ point
in the passing subset sits at $\Ne = 85.5$ ($\tau_* = 0.7996$,
$\Mphi/\Mpl = 41.07$).

The geometric reason is that the hilltop ($\eps=0$) is a slow-roll
\emph{attractor} on both sides. Departing the hilltop requires
$\delta\tau$ such that $\eps(\delta\tau) = 1$, but at
$\Mphi \approx 41\Mpl$ the corresponding $\delta\tau$ exceeds the
half-width of the available $\tau \in (1/\sqrt 5,\,1)$ window. The
inflaton effectively cannot reach $\eps = 1$ within the natural
parameter range; it slow-rolls to $\tau \to 1$ (vestigial melt)
without breaking slow-roll cleanly.

\section{Joint structural no-go}
\label{sec:joint}

Combining Sections~\ref{sec:eta} and~\ref{sec:efolds}:
\begin{enumerate}
\item Sub-Planckian inflaton scale ($\Mphi \le \Mpl$): $\nS$ misses the
   Planck $2\sigma$ window by orders of magnitude (the $\eta$-problem).
\item Super-Planckian inflaton scale ($\Mphi \in [37.1, 47.4]\Mpl$):
   $(\nS,\rT)$ lands inside the Planck/BICEP region but $\Ne$
   overshoots $[50,65]$.
\end{enumerate}
There is no choice of $(\fzero, \Mphi)$ for which the natural-branch
vestigial potential simultaneously satisfies Planck-$2\sigma$ on $\nS$
and BICEP/Keck-$95\%$ on $\rT$ and the canonical e-fold count.
Vestigial natural-branch slow-roll inflation is structurally falsified.

\paragraph{Caveat: non-natural branches.} We have only constrained the
\emph{natural} branch $\tau\in(1/\sqrt{5},\,1)$ of the vestigial
potential. Non-natural mechanisms — e.g., off-equilibrium melt with
back-reaction, multi-field uplift of $\tau$ to a positive-energy
configuration, or coupling to an external field that re-shapes the
effective potential — fall outside the scope of \texttt{VestigialEOS.lean}
and remain open. The structural no-go applies to vestigial inflation
\emph{as currently formalized}; not to all possible
vestigial-derived inflationary scenarios.

\section{Lean module summary}
\label{sec:lean}

The wave ships \texttt{lean/SKEFTHawking/VestigialInflationNoGo.lean}
with $15$ substantive theorems plus $1$ module summary marker, $0$
sorry, $0$ new axioms (verified \texttt{propext, Classical.choice,
Quot.sound} only on the load-bearing
\texttt{vestigial\_natural\_branch\_inflation\_falsified} and
\texttt{nSAtHilltop\_planck\_compatible\_implies\_super\_Planckian}
via \texttt{\#print axioms}).

The headline result theorems are:
\begin{itemize}
\item \texttt{etaAtHilltop\_eq\_neg\_thirty\_ratio\_sq} — closed form
   $\eta = -30(\Mpl/\Mphi)^2$.
\item \texttt{etaAtHilltop\_abs\_lt\_one\_implies\_super\_Planckian} —
   slow-roll viability requires $\Mphi^2 > 30\,\Mpl^2$.
\item \texttt{nSAtHilltop\_planck\_compatible\_implies\_super\_Planckian} —
   Planck $2\sigma$ on $\nS$ requires $\Mphi^2 > 1300\,\Mpl^2$.
\item \texttt{vestigial\_natural\_branch\_inflation\_falsified} —
   contrapositive: $\Mphi \le \Mpl \Longrightarrow \nS < 0.9565$.
\end{itemize}

The Python mirror at \texttt{src/vestigial\_inflation/} +
\texttt{tests/test\_vestigial\_inflation\_no\_go.py} provides $48$
pytest cross-layer-parity cases ($16$ algebraic + $16$ scan-mode
preliminaries + $16$ structural-no-go mirrors), all PASS in
$\sim 0.06$~s. The 2574-point $(\fzero,\Mphi,\tau_*)$ scan output is at
\texttt{figures/data/vestigial\_inflation\_preliminary\_scan.json}.

\section{Conclusion}

Vestigial natural-branch slow-roll inflation is structurally falsified
within the EFT validity window, by a closed-form geometric $\eta$-problem
combined with an e-fold overshoot. The result joins Phase~5y and
Phase~6b W2 in a unified ledger of vestigial natural-branch
falsifications: of the three cosmological observables to which the
framework can be applied (DESI dark energy, CMB-$\ell$ perturbations,
Planck/BICEP $(\nS,\rT)$), the natural branch fails all three. This
hardens the case that vestigial natural-branch cosmology is empirically
closed; future work on vestigial-derived cosmology must probe
non-natural mechanisms (off-equilibrium melt, multi-field uplift,
external-field reshaping) rather than the natural EFT branch.

\bibliographystyle{apsrev4-2}
\begin{thebibliography}{99}

\bibitem{Phase6bRoadmap}
SK-EFT Hawking Project Phase 6b Roadmap (2026-04-24, prepared from
Lit-Search/Phase-5z/Post-SK-EFT Research Program Strategy.md v2.0
\S5).

\bibitem{Planck2018Inflation}
Planck Collaboration (Y.~Akrami \emph{et~al.}),
\emph{Planck 2018 results.\ X.\ Constraints on inflation},
Astron.\ Astrophys.\ \textbf{641}, A10 (2020). DOI:
\href{https://doi.org/10.1051/0004-6361/201833887}{10.1051/0004-6361/201833887}.

\bibitem{BICEPKeck2021}
BICEP/Keck Collaboration (P.A.R.\ Ade \emph{et~al.}),
\emph{Improved constraints on primordial gravitational waves using
Planck, WMAP, and BICEP/Keck observations through the 2018 observing
season},
Phys.\ Rev.\ Lett.\ \textbf{127}, 151301 (2021). DOI:
\href{https://doi.org/10.1103/PhysRevLett.127.151301}{10.1103/PhysRevLett.127.151301}.

\bibitem{KoivistoNunesMulryne2012}
T.S.~Koivisto, D.~Nunes, D.~Mulryne,
\emph{The vector inflation versus the inflaton},
JCAP \textbf{1209}, 026 (2012); arXiv:1209.2156. (Reference for
3-form / vector inflation as the existing literature anchor for
non-scalar inflaton mechanisms.)

\end{thebibliography}

\end{document}
